\section*{अध्याय \ref{ch:b1:electrophilic-additions} --- ऐल्कीन और ऐल्काइन: इलेक्ट्रॉनरागी योगज}

\begin{solution}{exo:b1:electrophilic-additions:1}
2-क्लोरोप्रोपेन; 2-ब्रोमो-2-मेथिलप्रोपेन; 1-ब्रोमो-1-मेथिलसाइक्लोहेक्सेन;
2-क्लोरोब्यूट-2-ईन।
\end{solution}

\begin{solution}{exo:b1:electrophilic-additions:2}
1,2-डाइब्रोमोएथेन; \textit{trans}-1,2-डाइब्रोमोसाइक्लोहेक्सेन;
1,2-डाइब्रोमोप्रोपेन।
\end{solution}

\begin{solution}{exo:b1:electrophilic-additions:3}
प्रोपेन-2-ऑल; 2-मेथिलप्रोपेन-2-ऑल; 1-मेथिलसाइक्लोहेक्सेन-1-ऑल (हर स्थिति में
मार्कोवनिकोव)।
\end{solution}

\begin{solution}{exo:b1:electrophilic-additions:4}
प्रोपाइन (\omterm{def:b1:electrophilic-additions:acetylide}{अंतस्थ ऐल्काइन}): \ce{CH3C#CH + NH2- -> CH3C#C- + NH3}। एथेनॉल:
\ce{CH3CH2OH + NH2- -> CH3CH2O- + NH3}। ब्यूट-2-आइन और प्रोपीन में कोई
पर्याप्त अम्लीय हाइड्रोजन नहीं है।
\end{solution}

\begin{solution}{exo:b1:electrophilic-additions:5}
$\pi$ युग्म \ce{H+} को \ce{CH2} कार्बन पर ले लेता है: तृतीयक धनायन
\ce{(CH3)3C+}; जल का एक \omterm{def:b1:lewis-resonance-vsepr:lewis}{एकाकी युग्म} उससे आबंध बनाता है: \ce{(CH3)3COH2+}; एक जल
अणु एक प्रोटॉन लेता है: \ce{(CH3)3COH}। गरम सांद्र अम्ल में जल
दुर्लभ होता है और वाष्पशील ऐल्कीन निकल जाती है: उत्क्रम अभिक्रिया, \omterm{def:b1:alcohol-activation:dehydration}{निर्जलीकरण},
आगे बढ़ती है।
\end{solution}

\begin{solution}{exo:b1:electrophilic-additions:6}
2-मेथिलप्रोपीन का प्रोटॉनन तृतीयक धनायन देता है, प्रोपीन का
द्वितीयक। दर-निर्धारक प्रोटॉनन की \omterm{def:b1:elementary-steps:energy-profile}{संक्रमण अवस्था} धनायन से
मिलती-जुलती है (हैमंड): अधिक स्थायी तृतीयक धनायन तक
नीचे अवरोध से पहुँचा जाता है।
\end{solution}

\begin{solution}{exo:b1:electrophilic-additions:7}
ब्रोमोनियम आयन वलय के एक फलक पर सेतु बनाता है; ब्रोमाइड उसे
दूसरे फलक से खोलता है: दोनों ब्रोमीन \textit{trans} पहुँचते हैं। \textit{trans}
डाइब्रोमाइड \omterm{def:b1:stereochemistry-in-depth:stereocentre}{किरैल} है (कोई दर्पण-तल नहीं), पर ब्रोमाइड दोनों कार्बनों पर समान रूप से आक्रमण
करता है: दोनों \omterm{def:b1:stereochemistry-in-depth:enantiomers}{प्रतिबिंबरूपी} समान मात्रा में बनते हैं, और उत्पाद
रेसेमिक, प्रकाशतः निष्क्रिय है।
\end{solution}

\begin{solution}{exo:b1:electrophilic-additions:8}
C1 का प्रोटॉनन C2 पर द्वितीयक धनायन देता है, जो हाइड्रोजन वहन करने वाले तृतीयक
C3 के बगल में है; वह हाइड्रोजन अपने युग्म के साथ C2 पर खिसक जाता है (एक
1,2-हाइड्राइड विस्थापन), जिससे C3 पर एक तृतीयक धनायन बचता है, जिसे
\ce{Cl-} ग्रहण करता है: 2-क्लोरो-2-मेथिलब्यूटेन।
\end{solution}

\begin{solution}{exo:b1:electrophilic-additions:9}
प्रोपाइन + सोडियम ऐमाइड: प्रोपाइनाइड आयन; 1-ब्रोमोप्रोपेन के साथ (SN2):
\ce{CH3C#C-CH2CH2CH3}, हेक्स-2-आइन। बेंज़ैल्डिहाइड के साथ प्रोपाइनाइड, फिर तनु
अम्ल: \ce{C6H5CH(OH)C#CCH3}, 1-फ़ेनिलब्यूट-2-आइन-1-ऑल।
\end{solution}

\begin{solution}{exo:b1:electrophilic-additions:10}
\cip{E}-ब्यूट-2-ईन: एक फलक पर ब्रोमोनियम, दूसरे से C2 पर
या C3 पर ब्रोमाइड; दोनों आक्रमण वही यौगिक देते हैं, वर्णक \cip{2R,3S}:
\omterm{def:b1:stereochemistry-in-depth:enantiomers}{मेसो यौगिक}, एक \omterm{def:b1:stereochemistry-in-depth:isomers}{त्रिविम समावयवी}। \cip{Z}-ब्यूट-2-ईन: C2 पर आक्रमण
\cip{2S,3S} देता है (या दूसरे फलक से उसका दर्पण-प्रतिबिंब), C3 पर आक्रमण
\cip{2R,3R}; समान मात्राएँ: दो \omterm{def:b1:stereochemistry-in-depth:isomers}{त्रिविम समावयवी}, एक \omterm{def:b1:stereochemistry-in-depth:enantiomers}{रेसेमिक मिश्रण}।
\end{solution}

\begin{solution}{exo:b1:electrophilic-additions:11}
असममित ब्रोमोनियम आयन में अधिक प्रतिस्थापित कार्बन धन आवेश का
अधिक भाग वहन करता है (उसका C--Br आबंध लंबा और दुर्बल है, लगभग
तृतीयक या द्वितीयक धनायन जैसा)। जल, बड़े आधिक्य में, उसी कार्बन पर
सेतु के विपरीत फलक से आक्रमण करता है: 1-ब्रोमोप्रोपेन-2-ऑल।
\end{solution}

\begin{solution}{exo:b1:electrophilic-additions:12}
$D = 1$: एक ऐल्कीन। 2-मेथिलब्यूट-1-ईन और 2-मेथिलब्यूट-2-ईन दोनों
तृतीयक धनायन देती हैं, इसलिए 2-ब्रोमो-2-मेथिलब्यूटेन और 2-मेथिलब्यूटेन-2-ऑल। उनके
प्रोटॉन NMR स्पेक्ट्रम भिन्न हैं: पहली के लिए दो वाइनिलिक प्रोटॉन (\ce{=CH2}, एकक),
दूसरी के लिए एक (\ce{=CH-}, एक चतुष्क)।
\end{solution}

\begin{solution}{pb:b1:electrophilic-additions:1}
\textbf{1.} हर कार्बन पर \ce{CH3} का क्रम H से ऊपर है: \cip{Z} में दोनों मेथिल
एक ही ओर हैं, \cip{E} में विपरीत ओर।
\textbf{2.} ऐसे \omterm{def:b1:stereochemistry-in-depth:isomers}{त्रिविम समावयवी} जो दर्पण-प्रतिबिंब नहीं हैं: \omterm{def:b1:stereochemistry-in-depth:enantiomers}{अप्रतिबिंबरूपी}।
\textbf{3.} C=C के परितः घूर्णन $\pi$ आबंध को तोड़ देगा, जिसकी कीमत
कमरे के ताप पर संघट्टों द्वारा उपलब्ध ऊर्जा से कहीं अधिक है।
\textbf{4.} \cip{E}: उसके मेथिल समूह एक-दूसरे पर भीड़ नहीं करते।
\textbf{5.} \ce{C4H8 + Br2 -> C4H8Br2}।
\textbf{6.} दो तुल्य त्रिविम केंद्र: अधिक से अधिक चार, वास्तव में तीन
(\cip{2R,3R}, \cip{2S,3S} और मेसो \cip{2R,3S})।
\textbf{7.} $\pi$ युग्म \ce{Br2} के एक Br परमाणु पर आक्रमण करता है, Br--Br युग्म
\ce{Br-} के रूप में निकल जाता है, और आक्रांत ब्रोमीन का एक \omterm{def:b1:lewis-resonance-vsepr:lewis}{एकाकी युग्म}
दूसरे कार्बन से आबंध बनाता है: एक तीन-सदस्यीय ब्रोमोनियम आयन।
\textbf{8.} सेतुबद्ध आयन में हर परमाणु का अष्टक पूरा है; खुले द्वितीयक
\omterm{def:b1:electronic-effects:intermediates}{कार्बधनायन} में छह इलेक्ट्रॉनों वाला एक कार्बन होता।
\textbf{9.} \ce{Br-} का एक \omterm{def:b1:lewis-resonance-vsepr:lewis}{एकाकी युग्म} वलय के एक कार्बन से आबंध बनाता है; सेतु का C--Br आबंध
टूटता है, और उसका युग्म सेतु वाले ब्रोमीन पर चला जाता है।
\textbf{10.} सेतु एक फलक घेरता है; SN2 की तरह, \omterm{def:b1:electronic-effects:nucleophile}{नाभिकरागी}
टूटने वाले आबंध के विपरीत प्रवेश करता है।
\textbf{11.} एक \omterm{def:b1:electrophilic-additions:halonium}{प्रति-योगज}।
\textbf{12.} जल ब्रोमोनियम आयन के लिए ब्रोमाइड से प्रतिस्पर्धा करेगा और
एक ब्रोमोहाइड्रिन देगा।
\textbf{13.} \cip{2R,3S} (या \cip{2S,3R}, वही यौगिक)।
\textbf{14.} वही यौगिक।
\textbf{15.} ऐसे \omterm{def:b1:stereochemistry-in-depth:isomers}{संरूपण} में जिसमें दोनों ब्रोमीन ग्रसित हों, एक दर्पण-तल
C2 और C3 के बीच से गुज़रता है; वर्णक विपरीत हैं।
\textbf{16.} \cip{2R,3R} और \cip{2S,3S}।
\textbf{17.} समान मात्राएँ: एक \omterm{def:b1:stereochemistry-in-depth:enantiomers}{रेसेमिक मिश्रण}।
\textbf{18.} हाँ: दोनों \omterm{def:b1:stereochemistry-in-depth:enantiomers}{अप्रतिबिंबरूपी} ऐल्कीन भिन्न उत्पाद देती हैं
(मेसो बनाम रेसेमिक)।
\textbf{19.} कोई नहीं: दोनों \omterm{def:b1:stereochemistry-in-depth:enantiomers}{प्रतिबिंबरूपी} एक-दूसरे को निरस्त कर देते हैं।
\textbf{20.} हर एक $5.60/56.0 = \qty{0.100}{mol}$।
\textbf{21.} हर एक $0.100 \times 0.85 \times 215.8 = \qty{18.3}{g}$।
\textbf{22.} \textbf{$[\alpha] = 0^\circ$ (\omterm{def:b1:stereochemistry-in-depth:enantiomers}{मेसो यौगिक}), \qty{18.3}{g}}।
\end{solution}
